% --- Archon physics formalization source begin ---
% archon:physics
% archon:covers PhyXMiniProblems/problem_phyx_mini_0685.lean
% archon:source-report reports/phyx_mini/problem_phyx_mini_0685.source.json

\chapter{Physics problem phyx\_mini\_0685}
\label{ch:PhyXMiniProblems_problem_phyx_mini_0685}

\paragraph{Problem source.}
\#\# Physical scenario

A playground is on the flat roof of a city school, \textbackslash{}( 6.00 \textbackslash{}text\{ m\} \textbackslash{}) above the street below as shown in figure. The vertical wall of the building is \textbackslash{}( h = 7.00 \textbackslash{}text\{ m\} \textbackslash{}) high, forming a 1-m-high railing around the playground. A ball has fallen to the street below, and a passerby returns it by launching it at an angle of \textbackslash{}( \textbackslash{}theta = 53.0\textasciicircum{}\textbackslash{}circ \textbackslash{}) above the horizontal at a point \textbackslash{}( d = 24.0 \textbackslash{}text\{ m\} \textbackslash{}) from the base of the building wall. The ball takes \textbackslash{}( 2.20 \textbackslash{}text\{ s\} \textbackslash{}) to reach a point vertically above the wall.

\#\# Question

Find the horizontal distance from the wall to the point on the roof where the ball lands.

\#\# Answer choices

- A: A: \textbackslash{}( 2.86 \textbackslash{}text\{ m\} \textbackslash{})

- B: B: \textbackslash{}( 3.12 \textbackslash{}text\{ m\} \textbackslash{})

- C: C: \textbackslash{}( 2.79 \textbackslash{}text\{ m\} \textbackslash{})

- D: D: \textbackslash{}( 2.13 \textbackslash{}text\{ m\} \textbackslash{})

\#\# Figure caption (auxiliary; use the image as primary evidence)

The image depicts a scenario involving projectile motion. A person on the left is shown kicking a ball at an angle \textbackslash{}(\textbackslash{}theta\textbackslash{}) from the horizontal. The trajectory of the ball is shown as a dashed curve, demonstrating a typical parabolic path. 

To the right of the kicker, there is a building with two standing figures on its roof. The height of the building is labeled as \textbackslash{}(h\textbackslash{}), while the horizontal distance from the kicker to the building's wall is labeled as \textbackslash{}(d\textbackslash{}). 

The image is likely part of a physics problem illustrating the concepts of projectile motion, such as launch angle, initial velocity, range, and obstacle height.

\#\# Dataset metadata

Kinematics; Physical Model Grounding Reasoning, Spatial Relation Reasoning

\paragraph{Recorded answer/context.}
C: C: \textbackslash{}( 2.79 \textbackslash{}text\{ m\} \textbackslash{})

\paragraph{Figure/image path.}
/root/proposal\_for\_physic/science-mango/phyx\_mini\_run/phyx\_data/test\_image/685.png

\paragraph{Formalization target.}
create a compiling Lean file with sorry bodies at `PhyXMiniProblems/problem\_phyx\_mini\_0685.lean`. The Lean declarations must preserve the physical quantities, dimensions or dimensional roles, figure labels, governing-law hypotheses, and final relation expressed by this problem.
Use Mathlib/Physlib names found through LeanExplore where available. If a physics API is missing, introduce faithful local abstractions rather than scalar placeholder aliases.

\begin{theorem}[Physics formalization target]
\label{thm:physics:phyx_mini_0685:target}
\lean{PhyXMiniProblems.ProblemPhyXMini0685.problem_phyx_mini_0685}
\uses{def:physics:phyx-mini-0685:phyxminiproblems-problemphyxmini0685-lengthinmeters, def:physics:phyx-mini-0685:phyxminiproblems-problemphyxmini0685-rooftopprojectilesetup, def:physics:phyx-mini-0685:phyxminiproblems-problemphyxmini0685-matchesprimaryfigure, def:physics:phyx-mini-0685:phyxminiproblems-problemphyxmini0685-matchesproblemdescription, def:physics:phyx-mini-0685:phyxminiproblems-problemphyxmini0685-usesstandardnearearthgravity, def:physics:phyx-mini-0685:phyxminiproblems-problemphyxmini0685-hasphysicalprojectileparameters, def:physics:phyx-mini-0685:phyxminiproblems-problemphyxmini0685-satisfiesidealprojectilelaws, def:physics:phyx-mini-0685:phyxminiproblems-problemphyxmini0685-computedlandingdistanceinmeters, def:physics:phyx-mini-0685:phyxminiproblems-problemphyxmini0685-matchesdisplayeddistance, def:physics:phyx-mini-0685:phyxminiproblems-problemphyxmini0685-isuniqueclosestdisplayeddistance}
The assigned autoformalize agent should translate this physics problem into Lean declarations in the covered file, with theorem and lemma proof bodies written as `by sorry`.
\end{theorem}
\begin{proof}
This is an autoformalization task, not a proof task. Produce statements that can later be proved by the physics prover without weakening the physical model.
\end{proof}

% --- Archon declaration topology begin ---
\section{Declaration topology}
\begin{definition}[speed Dimension]
  \label{def:physics:phyx-mini-0685:phyxminiproblems-problemphyxmini0685-speeddimension}
  \lean{PhyXMiniProblems.ProblemPhyXMini0685.speedDimension}
  The physical dimension of speed, L T\textasciicircum{}-1
\end{definition}
\begin{proof}
  This is the declared model definition of speed Dimension.
\end{proof}
\begin{definition}[acceleration Dimension]
  \label{def:physics:phyx-mini-0685:phyxminiproblems-problemphyxmini0685-accelerationdimension}
  \lean{PhyXMiniProblems.ProblemPhyXMini0685.accelerationDimension}
  The physical dimension of acceleration, L T\textasciicircum{}-2
\end{definition}
\begin{proof}
  This is the declared model definition of acceleration Dimension.
\end{proof}
\begin{definition}[Length Quantity]
  \label{def:physics:phyx-mini-0685:phyxminiproblems-problemphyxmini0685-lengthquantity}
  \lean{PhyXMiniProblems.ProblemPhyXMini0685.LengthQuantity}
  A nonnegative physical length
\end{definition}
\begin{proof}
  This is the declared model definition of Length Quantity.
\end{proof}
\begin{definition}[Signed Length Quantity]
  \label{def:physics:phyx-mini-0685:phyxminiproblems-problemphyxmini0685-signedlengthquantity}
  \lean{PhyXMiniProblems.ProblemPhyXMini0685.SignedLengthQuantity}
  A signed one-dimensional physical displacement or coordinate
\end{definition}
\begin{proof}
  This is the declared model definition of Signed Length Quantity.
\end{proof}
\begin{definition}[Time Quantity]
  \label{def:physics:phyx-mini-0685:phyxminiproblems-problemphyxmini0685-timequantity}
  \lean{PhyXMiniProblems.ProblemPhyXMini0685.TimeQuantity}
  A nonnegative physical duration measured from launch
\end{definition}
\begin{proof}
  This is the declared model definition of Time Quantity.
\end{proof}
\begin{definition}[Speed Quantity]
  \label{def:physics:phyx-mini-0685:phyxminiproblems-problemphyxmini0685-speedquantity}
  \lean{PhyXMiniProblems.ProblemPhyXMini0685.SpeedQuantity}
  \uses{def:physics:phyx-mini-0685:phyxminiproblems-problemphyxmini0685-speeddimension}
  A nonnegative launch-speed magnitude
\end{definition}
\begin{proof}
  This is the declared model definition of Speed Quantity.
\end{proof}
\begin{definition}[Acceleration Quantity]
  \label{def:physics:phyx-mini-0685:phyxminiproblems-problemphyxmini0685-accelerationquantity}
  \lean{PhyXMiniProblems.ProblemPhyXMini0685.AccelerationQuantity}
  \uses{def:physics:phyx-mini-0685:phyxminiproblems-problemphyxmini0685-accelerationdimension}
  A nonnegative magnitude of gravitational acceleration
\end{definition}
\begin{proof}
  This is the declared model definition of Acceleration Quantity.
\end{proof}
\begin{definition}[length In Meters]
  \label{def:physics:phyx-mini-0685:phyxminiproblems-problemphyxmini0685-lengthinmeters}
  \lean{PhyXMiniProblems.ProblemPhyXMini0685.lengthInMeters}
  \uses{def:physics:phyx-mini-0685:phyxminiproblems-problemphyxmini0685-lengthquantity}
  Metre readout of a nonnegative physical length
\end{definition}
\begin{proof}
  This is the declared model definition of length In Meters.
\end{proof}
\begin{definition}[signed Length In Meters]
  \label{def:physics:phyx-mini-0685:phyxminiproblems-problemphyxmini0685-signedlengthinmeters}
  \lean{PhyXMiniProblems.ProblemPhyXMini0685.signedLengthInMeters}
  \uses{def:physics:phyx-mini-0685:phyxminiproblems-problemphyxmini0685-signedlengthquantity}
  Metre readout of a signed displacement coordinate
\end{definition}
\begin{proof}
  This is the declared model definition of signed Length In Meters.
\end{proof}
\begin{definition}[time In Seconds]
  \label{def:physics:phyx-mini-0685:phyxminiproblems-problemphyxmini0685-timeinseconds}
  \lean{PhyXMiniProblems.ProblemPhyXMini0685.timeInSeconds}
  \uses{def:physics:phyx-mini-0685:phyxminiproblems-problemphyxmini0685-timequantity}
  Second readout of a physical duration
\end{definition}
\begin{proof}
  This is the declared model definition of time In Seconds.
\end{proof}
\begin{definition}[speed In Meters Per Second]
  \label{def:physics:phyx-mini-0685:phyxminiproblems-problemphyxmini0685-speedinmeterspersecond}
  \lean{PhyXMiniProblems.ProblemPhyXMini0685.speedInMetersPerSecond}
  \uses{def:physics:phyx-mini-0685:phyxminiproblems-problemphyxmini0685-speedquantity}
  Metre-per-second readout of a speed magnitude
\end{definition}
\begin{proof}
  This is the declared model definition of speed In Meters Per Second.
\end{proof}
\begin{definition}[acceleration In Meters Per Second Squared]
  \label{def:physics:phyx-mini-0685:phyxminiproblems-problemphyxmini0685-accelerationinmeterspersecondsquared}
  \lean{PhyXMiniProblems.ProblemPhyXMini0685.accelerationInMetersPerSecondSquared}
  \uses{def:physics:phyx-mini-0685:phyxminiproblems-problemphyxmini0685-accelerationquantity}
  Metre-per-second-squared readout of an acceleration magnitude
\end{definition}
\begin{proof}
  This is the declared model definition of acceleration In Meters Per Second Squared.
\end{proof}
\begin{definition}[Projectile Model]
  \label{def:physics:phyx-mini-0685:phyxminiproblems-problemphyxmini0685-projectilemodel}
  \lean{PhyXMiniProblems.ProblemPhyXMini0685.ProjectileModel}
  The idealized environment in which the standard projectile equations hold
\end{definition}
\begin{proof}
  This is the declared model definition of Projectile Model.
\end{proof}
\begin{definition}[Launch Location]
  \label{def:physics:phyx-mini-0685:phyxminiproblems-problemphyxmini0685-launchlocation}
  \lean{PhyXMiniProblems.ProblemPhyXMini0685.LaunchLocation}
  The physical place from which the ball is returned
\end{definition}
\begin{proof}
  This is the declared model definition of Launch Location.
\end{proof}
\begin{definition}[Landing Surface]
  \label{def:physics:phyx-mini-0685:phyxminiproblems-problemphyxmini0685-landingsurface}
  \lean{PhyXMiniProblems.ProblemPhyXMini0685.LandingSurface}
  The surface on which the question says the ball later lands
\end{definition}
\begin{proof}
  This is the declared model definition of Landing Surface.
\end{proof}
\begin{definition}[Figure Object]
  \label{def:physics:phyx-mini-0685:phyxminiproblems-problemphyxmini0685-figureobject}
  \lean{PhyXMiniProblems.ProblemPhyXMini0685.FigureObject}
  Physical objects visibly represented in image 685.png
\end{definition}
\begin{proof}
  This is the declared model definition of Figure Object.
\end{proof}
\begin{definition}[Figure Label]
  \label{def:physics:phyx-mini-0685:phyxminiproblems-problemphyxmini0685-figurelabel}
  \lean{PhyXMiniProblems.ProblemPhyXMini0685.FigureLabel}
  Mathematical labels visibly printed in the primary image
\end{definition}
\begin{proof}
  This is the declared model definition of Figure Label.
\end{proof}
\begin{definition}[Figure Anchor]
  \label{def:physics:phyx-mini-0685:phyxminiproblems-problemphyxmini0685-figureanchor}
  \lean{PhyXMiniProblems.ProblemPhyXMini0685.FigureAnchor}
  Schematic points needed to retain the spatial relations in the figure
\end{definition}
\begin{proof}
  This is the declared model definition of Figure Anchor.
\end{proof}
\begin{definition}[Rooftop Projectile Figure]
  \label{def:physics:phyx-mini-0685:phyxminiproblems-problemphyxmini0685-rooftopprojectilefigure}
  \lean{PhyXMiniProblems.ProblemPhyXMini0685.RooftopProjectileFigure}
  \uses{def:physics:phyx-mini-0685:phyxminiproblems-problemphyxmini0685-figureobject, def:physics:phyx-mini-0685:phyxminiproblems-problemphyxmini0685-figurelabel, def:physics:phyx-mini-0685:phyxminiproblems-problemphyxmini0685-figureanchor}
  Qualitative and schematic-coordinate evidence from the primary raster. Coordinates are dimensionless drawing coordinates, not measured lengths
\end{definition}
\begin{proof}
  This is the declared model definition of Rooftop Projectile Figure.
\end{proof}
\begin{definition}[Rooftop Projectile Setup]
  \label{def:physics:phyx-mini-0685:phyxminiproblems-problemphyxmini0685-rooftopprojectilesetup}
  \lean{PhyXMiniProblems.ProblemPhyXMini0685.RooftopProjectileSetup}
  \uses{def:physics:phyx-mini-0685:phyxminiproblems-problemphyxmini0685-lengthquantity, def:physics:phyx-mini-0685:phyxminiproblems-problemphyxmini0685-signedlengthquantity, def:physics:phyx-mini-0685:phyxminiproblems-problemphyxmini0685-timequantity, def:physics:phyx-mini-0685:phyxminiproblems-problemphyxmini0685-speedquantity, def:physics:phyx-mini-0685:phyxminiproblems-problemphyxmini0685-accelerationquantity, def:physics:phyx-mini-0685:phyxminiproblems-problemphyxmini0685-projectilemodel, def:physics:phyx-mini-0685:phyxminiproblems-problemphyxmini0685-launchlocation, def:physics:phyx-mini-0685:phyxminiproblems-problemphyxmini0685-landingsurface, def:physics:phyx-mini-0685:phyxminiproblems-problemphyxmini0685-rooftopprojectilefigure}
  The independent physical quantities in the projectile experiment. The two coordinate functions give signed displacements from the street-level launch point as functions of elapsed physical time. In particular, landingDistanceFromWall is not defined from an answer choice
\end{definition}
\begin{proof}
  This is the declared model definition of Rooftop Projectile Setup.
\end{proof}
\begin{definition}[Matches Primary Figure]
  \label{def:physics:phyx-mini-0685:phyxminiproblems-problemphyxmini0685-matchesprimaryfigure}
  \lean{PhyXMiniProblems.ProblemPhyXMini0685.MatchesPrimaryFigure}
  \uses{def:physics:phyx-mini-0685:phyxminiproblems-problemphyxmini0685-rooftopprojectilesetup}
  Every named object and label is shown. The schematic positions say that the launch occurs at street level to the left of the vertical wall, the roof lies above the street and to the right of the wall, and the railing top marked h lies above the roof. No landing-distance readout occurs in the image
\end{definition}
\begin{proof}
  This is the declared model definition of Matches Primary Figure.
\end{proof}
\begin{definition}[degrees]
  \label{def:physics:phyx-mini-0685:phyxminiproblems-problemphyxmini0685-degrees}
  \lean{PhyXMiniProblems.ProblemPhyXMini0685.degrees}
  Convert a degree readout to a physical angle modulo 2 * pi
\end{definition}
\begin{proof}
  This is the declared model definition of degrees.
\end{proof}
\begin{definition}[Matches Problem Description]
  \label{def:physics:phyx-mini-0685:phyxminiproblems-problemphyxmini0685-matchesproblemdescription}
  \lean{PhyXMiniProblems.ProblemPhyXMini0685.MatchesProblemDescription}
  \uses{def:physics:phyx-mini-0685:phyxminiproblems-problemphyxmini0685-lengthinmeters, def:physics:phyx-mini-0685:phyxminiproblems-problemphyxmini0685-timeinseconds, def:physics:phyx-mini-0685:phyxminiproblems-problemphyxmini0685-rooftopprojectilesetup, def:physics:phyx-mini-0685:phyxminiproblems-problemphyxmini0685-degrees}
  The qualitative scenario and numerical quantities stated in the prose
\end{definition}
\begin{proof}
  This is the declared model definition of Matches Problem Description.
\end{proof}
\begin{definition}[Uses Standard Near Earth Gravity]
  \label{def:physics:phyx-mini-0685:phyxminiproblems-problemphyxmini0685-usesstandardnearearthgravity}
  \lean{PhyXMiniProblems.ProblemPhyXMini0685.UsesStandardNearEarthGravity}
  \uses{def:physics:phyx-mini-0685:phyxminiproblems-problemphyxmini0685-accelerationinmeterspersecondsquared, def:physics:phyx-mini-0685:phyxminiproblems-problemphyxmini0685-rooftopprojectilesetup}
  The standard near-Earth value needed by the textbook projectile model. It is a physical calibration, not a landing-distance assumption
\end{definition}
\begin{proof}
  This is the declared model definition of Uses Standard Near Earth Gravity.
\end{proof}
\begin{definition}[Has Physical Projectile Parameters]
  \label{def:physics:phyx-mini-0685:phyxminiproblems-problemphyxmini0685-hasphysicalprojectileparameters}
  \lean{PhyXMiniProblems.ProblemPhyXMini0685.HasPhysicalProjectileParameters}
  \uses{def:physics:phyx-mini-0685:phyxminiproblems-problemphyxmini0685-lengthinmeters, def:physics:phyx-mini-0685:phyxminiproblems-problemphyxmini0685-timeinseconds, def:physics:phyx-mini-0685:phyxminiproblems-problemphyxmini0685-speedinmeterspersecond, def:physics:phyx-mini-0685:phyxminiproblems-problemphyxmini0685-accelerationinmeterspersecondsquared, def:physics:phyx-mini-0685:phyxminiproblems-problemphyxmini0685-rooftopprojectilesetup}
  Positivity, ordering, and the upward/rightward branch depicted in the figure. These conditions do not assign a numerical value to the landing distance
\end{definition}
\begin{proof}
  This is the declared model definition of Has Physical Projectile Parameters.
\end{proof}
\begin{definition}[Satisfies Ideal Projectile Laws]
  \label{def:physics:phyx-mini-0685:phyxminiproblems-problemphyxmini0685-satisfiesidealprojectilelaws}
  \lean{PhyXMiniProblems.ProblemPhyXMini0685.SatisfiesIdealProjectileLaws}
  \uses{def:physics:phyx-mini-0685:phyxminiproblems-problemphyxmini0685-timequantity, def:physics:phyx-mini-0685:phyxminiproblems-problemphyxmini0685-lengthinmeters, def:physics:phyx-mini-0685:phyxminiproblems-problemphyxmini0685-signedlengthinmeters, def:physics:phyx-mini-0685:phyxminiproblems-problemphyxmini0685-timeinseconds, def:physics:phyx-mini-0685:phyxminiproblems-problemphyxmini0685-speedinmeterspersecond, def:physics:phyx-mini-0685:phyxminiproblems-problemphyxmini0685-accelerationinmeterspersecondsquared, def:physics:phyx-mini-0685:phyxminiproblems-problemphyxmini0685-rooftopprojectilesetup}
  The governing ideal-projectile relations in coherent SI readouts: horizontal displacement is uniform motion with velocity v0 * cos(theta); vertical displacement is v0 * sin(theta) * t - g * t\textasciicircum{}2 / 2; at the observed 2.20 s, the ball is vertically above the wall and clears the 7.00 m railing top; at the later landing time, the vertical coordinate is the 6.00 m roof height; and the landing distance is the horizontal displacement beyond the wall. These are general physical and geometric laws.
\end{definition}
\begin{proof}
  This is the declared model definition of Satisfies Ideal Projectile Laws.
\end{proof}
\begin{definition}[Answer Choice]
  \label{def:physics:phyx-mini-0685:phyxminiproblems-problemphyxmini0685-answerchoice}
  \lean{PhyXMiniProblems.ProblemPhyXMini0685.AnswerChoice}
  The four answer labels printed with the problem
\end{definition}
\begin{proof}
  This is the declared model definition of Answer Choice.
\end{proof}
\begin{definition}[answer Distance In Meters]
  \label{def:physics:phyx-mini-0685:phyxminiproblems-problemphyxmini0685-answerdistanceinmeters}
  \lean{PhyXMiniProblems.ProblemPhyXMini0685.answerDistanceInMeters}
  \uses{def:physics:phyx-mini-0685:phyxminiproblems-problemphyxmini0685-answerchoice}
  Distance in metres printed beside each displayed answer label
\end{definition}
\begin{proof}
  This is the declared model definition of answer Distance In Meters.
\end{proof}
\begin{definition}[recorded Answer Choice]
  \label{def:physics:phyx-mini-0685:phyxminiproblems-problemphyxmini0685-recordedanswerchoice}
  \lean{PhyXMiniProblems.ProblemPhyXMini0685.recordedAnswerChoice}
  \uses{def:physics:phyx-mini-0685:phyxminiproblems-problemphyxmini0685-answerchoice}
  The answer label recorded in the supplied dataset metadata
\end{definition}
\begin{proof}
  This is the declared model definition of recorded Answer Choice.
\end{proof}
\begin{definition}[computed Landing Time In Seconds]
  \label{def:physics:phyx-mini-0685:phyxminiproblems-problemphyxmini0685-computedlandingtimeinseconds}
  \lean{PhyXMiniProblems.ProblemPhyXMini0685.computedLandingTimeInSeconds}
  \uses{def:physics:phyx-mini-0685:phyxminiproblems-problemphyxmini0685-degrees}
  The later root of the roof-height quadratic, after eliminating the launch speed using the measured horizontal wall-crossing time
\end{definition}
\begin{proof}
  This is the declared model definition of computed Landing Time In Seconds.
\end{proof}
\begin{definition}[computed Landing Distance In Meters]
  \label{def:physics:phyx-mini-0685:phyxminiproblems-problemphyxmini0685-computedlandingdistanceinmeters}
  \lean{PhyXMiniProblems.ProblemPhyXMini0685.computedLandingDistanceInMeters}
  \uses{def:physics:phyx-mini-0685:phyxminiproblems-problemphyxmini0685-computedlandingtimeinseconds}
  The exact unrounded horizontal distance beyond the wall predicted by the stated data and the standard-gravity projectile model. The independent setup field is not definitionally equal to this expression
\end{definition}
\begin{proof}
  This is the declared model definition of computed Landing Distance In Meters.
\end{proof}
\begin{definition}[Matches Displayed Distance]
  \label{def:physics:phyx-mini-0685:phyxminiproblems-problemphyxmini0685-matchesdisplayeddistance}
  \lean{PhyXMiniProblems.ProblemPhyXMini0685.MatchesDisplayedDistance}
  \uses{def:physics:phyx-mini-0685:phyxminiproblems-problemphyxmini0685-lengthquantity, def:physics:phyx-mini-0685:phyxminiproblems-problemphyxmini0685-lengthinmeters, def:physics:phyx-mini-0685:phyxminiproblems-problemphyxmini0685-answerchoice, def:physics:phyx-mini-0685:phyxminiproblems-problemphyxmini0685-answerdistanceinmeters}
  Agreement with a displayed distance rounded to the nearest centimetre
\end{definition}
\begin{proof}
  This is the declared model definition of Matches Displayed Distance.
\end{proof}
\begin{definition}[Is Closest Displayed Distance]
  \label{def:physics:phyx-mini-0685:phyxminiproblems-problemphyxmini0685-isclosestdisplayeddistance}
  \lean{PhyXMiniProblems.ProblemPhyXMini0685.IsClosestDisplayedDistance}
  \uses{def:physics:phyx-mini-0685:phyxminiproblems-problemphyxmini0685-lengthquantity, def:physics:phyx-mini-0685:phyxminiproblems-problemphyxmini0685-lengthinmeters, def:physics:phyx-mini-0685:phyxminiproblems-problemphyxmini0685-answerchoice, def:physics:phyx-mini-0685:phyxminiproblems-problemphyxmini0685-answerdistanceinmeters}
  A displayed choice is at least as close as every other choice
\end{definition}
\begin{proof}
  This is the declared model definition of Is Closest Displayed Distance.
\end{proof}
\begin{definition}[Is Unique Closest Displayed Distance]
  \label{def:physics:phyx-mini-0685:phyxminiproblems-problemphyxmini0685-isuniqueclosestdisplayeddistance}
  \lean{PhyXMiniProblems.ProblemPhyXMini0685.IsUniqueClosestDisplayedDistance}
  \uses{def:physics:phyx-mini-0685:phyxminiproblems-problemphyxmini0685-lengthquantity, def:physics:phyx-mini-0685:phyxminiproblems-problemphyxmini0685-answerchoice, def:physics:phyx-mini-0685:phyxminiproblems-problemphyxmini0685-isclosestdisplayeddistance}
  A displayed choice is the unique closest printed value
\end{definition}
\begin{proof}
  This is the declared model definition of Is Unique Closest Displayed Distance.
\end{proof}
\begin{lemma}[landing Distance eq computed Expression]
  \label{lem:physics:phyx-mini-0685:phyxminiproblems-problemphyxmini0685-landingdistance-eq-computedexpression}
  \lean{PhyXMiniProblems.ProblemPhyXMini0685.landingDistance_eq_computedExpression}
  \uses{def:physics:phyx-mini-0685:phyxminiproblems-problemphyxmini0685-lengthinmeters, def:physics:phyx-mini-0685:phyxminiproblems-problemphyxmini0685-rooftopprojectilesetup, def:physics:phyx-mini-0685:phyxminiproblems-problemphyxmini0685-matchesproblemdescription, def:physics:phyx-mini-0685:phyxminiproblems-problemphyxmini0685-usesstandardnearearthgravity, def:physics:phyx-mini-0685:phyxminiproblems-problemphyxmini0685-hasphysicalprojectileparameters, def:physics:phyx-mini-0685:phyxminiproblems-problemphyxmini0685-satisfiesidealprojectilelaws, def:physics:phyx-mini-0685:phyxminiproblems-problemphyxmini0685-computedlandingdistanceinmeters}
  Eliminating the launch speed and selecting the later roof intersection gives the exact expression above. This is a derived conclusion rather than a field of the physical-law premise
\end{lemma}
\begin{proof}
  The conclusion follows from the governing relations and figure readouts named in the statement of landing Distance eq computed Expression.
\end{proof}
% --- Archon declaration topology end ---
% --- Archon physics formalization source end ---
